
\documentclass[11pt,twoside]{article}

% Do NOT use ANY packages other than asp2023.
\usepackage{asp2023}

\aspSuppressVolSlug
\resetcounters

\bibliographystyle{asp2021}

\markboth{Valame}{Seeing the Same Sky}

\begin{document}

\title{Seeing the Same Sky: Visual Thinking as a Collaboration Tool in the SKAO}

\author{Snehal Valame}

\affil{Sanikaizen Solutions, Pune, India; \email{snehal.valame@gmail.com}}

\paperauthor{Snehal Valame}{snehal.valame@gmail.com}{}{Sanikaizen Solutions}{ }{Pune}{Maharashtra}{ }{India}

\begin{abstract}
Large, geographically distributed software programmes bring together
stakeholders with different roles and mental models. Abundant
requirements, diagrams, and tracking tools do not guarantee shared
understanding. This paper describes the use of simple hand-drawn
visualisations as a collaboration tool within the Square Kilometre
Array Observatory (SKAO) Observation Management and Controls
software programme. Sketches supported conversations about
architecture, goals, capacity, delivery risks, and non-functional
requirements. By making systems, dependencies, and trade-offs
visible, they provided a common reference for discussion.
Visual thinking is not about polished diagrams, but enabling
meaningful conversations and shared understanding.
\end{abstract}


\section{Introduction}

In SKAO's distributed software programme, people may work
towards the same goal while seeing it differently: a scientist
sees a discovery, an operator a telescope to manage, an engineer
a potential defect, and a Product Owner a milestone to deliver.

As a Release Train Engineer within the Observation Management
and Controls (OMC) group, I have found hand-drawn visuals
useful for making mental models visible and creating a common
point of reference. These sketches complement formal diagrams
and documentation by helping people see, discuss, and refine ideas.

\articlefigure[width=0.60\textwidth]{C601_f1.eps}{fig:sky}
{Different stakeholders view the same system differently.}


\section{From Information to Shared Understanding}

Requirements, diagrams, Jira tickets, and discussions provide
information, but not necessarily shared understanding. A
networking handshake (SYN, SYN-ACK, ACK) establishes a
connection through messages and acknowledgements. Human
collaboration is less deterministic, particularly across time
zones: sending information does not guarantee understanding.

Hand-drawn visuals externalise assumptions and mental models,
giving people something concrete to question and correct.
The drawing is not the outcome; the conversation is.

\articlefigure[width=0.42\textwidth]{C601_f2.eps}{fig:communication}
{A visual analogy for communication and acknowledgement.}


\section{Making the Science Journey Visible}

The science journey is simple: plan an observation, operate
the telescope, process signals, and deliver science data.
Behind it lies a complex software ecosystem. Components may
progress independently, but end-to-end readiness depends on
their integration. A delay in telescope control or downstream
processing can affect the complete science journey.
A workflow sketch shifts discussion from individual components
to their dependencies, connections, and integration gaps.

\articlefigure[width=0.70\textwidth]{C601_f3.eps}{fig:workflow}
{The science journey and its software dependencies.}


\section{Visualising Goals and Impacts}

Impact Mapping starts with the goal, identifies actors, explores
the behavioural changes needed, and then considers deliverables.
For SKA-Low, the goal was to establish the route to Array
Assembly 2 (AA2) commissioning and science verification.
Commissioning scientists and telescope operators helped identify
the impacts and supporting capabilities, such as proposal
submission, user registration, and observation monitoring.
This shifts the conversation from what teams should build to
what is needed to achieve the goal.

\articlefigure[width=0.60\textwidth]{C601_f4.eps}{fig:impact}
{Connecting goals, actors, impacts, and deliverables.}


\section{Making Capacity Visible}

Engineering capacity is finite. Teams balance new capabilities
with reliability, defects, technical debt, infrastructure,
testing, and risks. Visualising capacity makes competing
investments tangible and supports prioritisation.

Not all capacity can be committed to planned features; some
must be reserved for interruptions. The discussion becomes:
what MUST, SHOULD, COULD, or WON'T be delivered?

\articlefigure[width=0.30\textwidth]{C601_f5.eps}{fig:capacity}
{Capacity allocation and competing engineering investments.}


\section{Making Delivery Risks Visible}

A critical defect can become a roadblock. When a showstopper
prevents a capability from operating, planned work may stop.
Visualising defects helps discuss severity, age, release impact,
and the effort required to resolve them. Reserved capacity
helps address such issues and resume planned delivery.

\articlefigure[width=0.55\textwidth]{C601_f7.eps}{fig:roadblock}
{A critical defect blocking planned delivery.}


\section{Looking Beyond the Demo}

A successful demo does not necessarily mean software is ready
for operations. Recoverability, maintainability, documentation,
interoperability, and unhappy-path scenarios matter when
software is integrated, operated, or handed over.

Visualising non-functional requirements (NFRs) brings these
qualities into scope discussions. They belong in the backlog
and Definition of Done. If future teams cannot use, test,
maintain, and evolve the software, the work is not complete.

\articlefigure[width=0.55\textwidth]{C601_f6.eps}{fig:nfr}
{NFRs support operational readiness and maintainability.}


\section{Reflection and Conclusion}

These sketches make the invisible visible: systems, workflows,
goals, dependencies, risks, and trade-offs. They create a
common reference for challenging assumptions and connecting
different responsibilities to one scientific journey.

Visual thinking does not replace communication; it gives
people something concrete to communicate about.

\end{document}
